\begin{table}[H]
  \centering
  \caption{Definitions of the baselines. All constrained baselines use the same
  candidate action set, feasibility check, channel warm-up model, and six-time-step
  minimum dwell-time constraint as PD-PPO unless
  stated otherwise.}
  \label{tab:executable_baseline_definitions}
  \scriptsize
  \renewcommand{\arraystretch}{0.92}
  \setlength{\tabcolsep}{2.6pt}
  \begin{tabular}{@{}p{0.19\linewidth}p{0.31\linewidth}p{0.25\linewidth}p{0.19\linewidth}@{}}
    \toprule
    Baseline & Decision rule & Selection and tie break & Information availability \\
    \midrule
    Fixed-schedule baseline
      & Apply one fixed channel subset.
      & Evaluate all six subsets on calibration/validation windows; minimize the
        macro-averaged normalized forecast loss, then mean forecast loss.
      & Feasible baseline; no condition labels at decision time. \\
    Fixed-priority baseline
      & Assign linearly decreasing scores by channel index at every time step.
      & Deterministic projector order.
      & Online and non-learned. \\
    Round-robin baseline
      & Give the highest scores to a cyclic group of two channels.
      & The time step and channel index define a fixed cyclic schedule; projection resolves conflicts.
      & Online and non-learned. \\
    AoI-priority baseline
      & Score each channel by clipped time since its last ready sampling opportunity.
      & Project the descending score order.
      & Online and non-learned. \\
    Random-feasible baseline
      & Draw independent Gaussian channel scores at each time step.
      & One random stream per simulation seed, seeded by simulation seed plus 100; the same 512-score stream restarts at each event-balanced test-window reset; project the sampled ranking.
      & Online stochastic baseline; no within-seed Monte Carlo averaging. \\
    Double DQN training configuration
      & Evaluate six Q-values and greedily select the largest currently feasible
        action during test evaluation.
      & Double DQN target with a feasible-action mask; deterministic
        lowest-index tie break.
      & Online learned baseline; no condition labels at decision time. \\
    Warning-rule baseline
      & Choose the largest particle, flux, or thermal warning signal when it is
        at least 0.5; otherwise choose the calm action.
      & On calibration/validation windows, separately minimize calm or event-type loss for
        each warning-signal case; break ties by mean forecast loss.
      & Uses simulated warning signals; no condition labels at decision time. \\
    Privileged one-step look-ahead reference using next-step test targets
      & Simulate every candidate for one time step and select the smallest next-step
        forecast loss on the test partition under the frozen primary forecaster.
      & Lower switching fraction, then lower action index.
      & Uses next-step test targets unavailable to the deployed scheduler; the objective omits minimum dwell-time, warm-up, and observation-age consequences. \\
    Privileged true-label baseline
      & Choose the optional channel selected on the calibration/validation partition using the exact
        label sequence $c_t,\ldots,c_{t+16}$.
      & Use the first active event type in the future-label interval; break calibration/validation
        ties by mean forecast loss.
      & Uses condition labels unavailable to the deployed scheduler. \\
    Full-observation baseline
      & Request all six channels at every time step.
      & No selection.
      & Deliberately removes the constraints; unconstrained benchmark. \\
    \bottomrule
  \end{tabular}
\end{table}
